\documentclass[12pt,a4paper]{article}
\usepackage[T2A]{fontenc}
\usepackage[utf8]{inputenc}
\usepackage[russian]{babel}
\usepackage{amsmath,amssymb,geometry,array,hyperref}
\geometry{margin=18mm}
\hypersetup{hidelinks}
\setlength{\parindent}{0pt}
\setlength{\parskip}{5pt}
\newcommand{\head}[1]{\medskip\textbf{#1}\par}
\begin{document}
\begin{center}\large\textbf{Лабораторная работа №4. Динамическое программирование}\end{center}
Борисюк Е. Г. Вариант 4.

\head{1. Условие и математическая постановка}
Компания распределяет 7 одинаковых серверных модулей между тремя сервисами.
Производительность сервисов складывается. Сервис обработки заказов должен
получить не менее 1 модуля; ни один сервис не может получить более 3 модулей.
Требуется распределить все 7 модулей так, чтобы получить максимальную
суммарную производительность.

\begin{center}\begin{tabular}{c|r|r|r}
$u$ & Обработка заказов & Аналитика & Поиск\\\hline
0&0&0&0\\1&6&2&4\\2&10&7&9\\3&13&12&13\\
4&16&16&15\\5&18&19&17\\6&19&21&18\\7&20&23&19
\end{tabular}\end{center}
Пусть $x_1,x_2,x_3$ --- количество модулей у обработки заказов, аналитики
и поиска; $g_1,g_2,g_3$ --- табличные функции производительности.
\[Z=g_1(x_1)+g_2(x_2)+g_3(x_3)\longrightarrow\max.\]
\[x_1+x_2+x_3=7,\quad 1\le x_1\le3,\quad0\le x_2,x_3\le3,
\quad x_1,x_2,x_3\in\mathbb Z.\]
Строки $u>3$ сохраняются как часть условия, но исключены из допустимых управлений.

\head{2. Этапы, состояния и уравнение Беллмана}
Этап $k$ соответствует сервису $k$. Состояние $s$ --- число модулей,
которое ещё требуется полностью распределить между сервисами $k,\ldots,3$.
Управление $u$ --- число модулей текущему сервису; следующее состояние $s-u$.
$F_k(s)$ --- максимальная производительность сервисов $k,\ldots,3$
при точном распределении $s$ модулей.
\[A_1=\{1,2,3\},\qquad A_2=A_3=\{0,1,2,3\}.\]
\[F_k(s)=\max_{u\in A_k,\ u\le s}\{g_k(u)+F_{k+1}(s-u)\}.\]
Граница: $F_4(0)=0$, $F_4(s)=-\infty$ при $s>0$.
Недостижимому состоянию соответствует $-\infty$; максимум по пустому
множеству равен $-\infty$. Недостижимые продолжения исключаются.

\newpage
\head{3. Обратный проход: поиск и аналитика}
Условная оптимизация идёт от последнего сервиса к первому: $F_3,F_2,F_1$.
В $U_k(s)$ сохраняются все управления, достигающие максимума.

Этап 3: весь остаток получает поиск. При $0\le s\le3$ имеем
$F_3(s)=g_3(s)$; при $s\ge4$ состояние недостижимо.
\begin{center}\begin{tabular}{c|rrrrrrrr}
$s$&0&1&2&3&4&5&6&7\\\hline
$F_3(s)$&0&4&9&13&$-\infty$&$-\infty$&$-\infty$&$-\infty$\\
$U_3(s)$&$\{0\}$&$\{1\}$&$\{2\}$&$\{3\}$&$\varnothing$&$\varnothing$&$\varnothing$&$\varnothing$
\end{tabular}\end{center}

Этап 2: перебираем число модулей аналитики. Кандидаты
$g_2(u)+F_3(s-u)$ перечислены в порядке возрастания допустимого $u$;
недостижимые продолжения не включены.
\begin{center}\begin{tabular}{c|c|l|r|c}
$s$&Допустимые $u$&Кандидаты&$F_2(s)$&$U_2(s)$\\\hline
0&$\{0\}$&0&0&$\{0\}$\\
1&$\{0,1\}$&4; 2&4&$\{0\}$\\
2&$\{0,1,2\}$&9; 6; 7&9&$\{0\}$\\
3&$\{0,1,2,3\}$&13; 11; 11; 12&13&$\{0\}$\\
4&$\{1,2,3\}$&15; 16; 16&16&$\{2,3\}$\\
5&$\{2,3\}$&20; 21&21&$\{3\}$\\
6&$\{3\}$&25&25&$\{3\}$\\
7&$\varnothing$&нет&$-\infty$&$\varnothing$
\end{tabular}\end{center}
Например, при $s=4$:
\begin{align*}
F_2(4)&=\max\{2+13;7+9;12+4\}\\
&=\max\{15;16;16\}=16,\quad U_2(4)=\{2,3\}.
\end{align*}
При $s=5$: $F_2(5)=\max\{7+13;12+9\}=21$, $U_2(5)=\{3\}$.
При $s=6$ допустимо только $u=3$: $F_2(6)=12+13=25$.
Распределить 7 модулей между двумя сервисами с верхней границей 3 невозможно.

\newpage
\head{4. Обратный проход: обработка заказов}
Используем найденные $F_2(s-u)$; управление $u\in\{1,2,3\}$.
Для $s=0$ допустимых управлений нет: $F_1(0)=-\infty$, $U_1(0)=\varnothing$.
\begin{center}\begin{tabular}{c|c|l|r|c}
$s$&Допустимые $u$&Кандидаты&$F_1(s)$&$U_1(s)$\\\hline
1&$\{1\}$&6&6&$\{1\}$\\
2&$\{1,2\}$&10; 10&10&$\{1,2\}$\\
3&$\{1,2,3\}$&15; 14; 13&15&$\{1\}$\\
4&$\{1,2,3\}$&19; 19; 17&19&$\{1,2\}$\\
5&$\{1,2,3\}$&22; 23; 22&23&$\{2\}$\\
6&$\{1,2,3\}$&27; 26; 26&27&$\{1\}$\\
7&$\{1,2,3\}$&31; 31; 29&31&$\{1,2\}$
\end{tabular}\end{center}
\begin{align*}
F_1(7)&=\max\{6+F_2(6);10+F_2(5);13+F_2(4)\}\\
&=\max\{6+25;10+21;13+16\}=31.
\end{align*}
$U_1(7)=\{1,2\}$, необходимо сохранить обе ветви.

\head{5. Прямой проход: восстановление стратегий}
Начинаем с $k=1$, $s=7$. Каждый раз заменяем $s$ на $s-u$.
\begin{center}\begin{tabular}{c|c|c|c|c}
Ветвь&Этап 1&Этап 2&Этап 3&План\\\hline
1&$s=7,u=1$&$s=6,u=3$&$s=3,u=3$&$(1,3,3)$\\
2&$s=7,u=2$&$s=5,u=3$&$s=2,u=2$&$(2,3,2)$
\end{tabular}\end{center}
Конечный остаток равен нулю. На этапах 2 и 3 в этих состояниях управление
единственно, поэтому других оптимальных стратегий нет.
Равенство в $U_2(4)$ относится к неоптимальной ветви $x_1=3$.

\head{6. Ответ и проверка}
\[\boxed{Z_{\max}=31,\quad(x_1,x_2,x_3)\in\{(1,3,3),(2,3,2)\}.}\]
\[g_1(1)+g_2(3)+g_3(3)=6+12+13=31,\]
\[g_1(2)+g_2(3)+g_3(2)=10+12+9=31.\]
В обоих случаях сумма модулей равна 7, все компоненты целые и не превышают 3;
обработка заказов получает хотя бы 1 модуль.

\newpage
\head{7. Псевдокод основного алгоритма}
$n$ --- число сервисов, $R$ --- общий ресурс, $A[k]$ --- допустимые управления.
\begin{tabbing}
\quad\=\quad\=\quad\=\quad\=\kill
Для всех $k=1,\ldots,n+1$ и $s=0,\ldots,R$: $F[k,s]\gets-\infty$\\
$F[n+1,0]\gets0$; все $U[k,s]\gets[\ ]$\\
Для $k$ от $n$ до 1:\\
\>Для $s$ от 0 до $R$:\\
\>\>Для $u\in A[k]$, где $u\le s$:\\
\>\>\>Если $F[k+1,s-u]=-\infty$: пропустить $u$\\
\>\>\>$v\gets g[k,u]+F[k+1,s-u]$\\
\>\>\>Если $v>F[k,s]$: $F[k,s]\gets v$; $U[k,s]\gets[u]$\\
\>\>\>Иначе если $v=F[k,s]$: добавить $u$ в $U[k,s]$\\
Восстановить$(k,s)$:\\
\>Если $k=n+1$: вернуть [пустой план] при $s=0$, иначе $[\ ]$\\
\>$plans\gets[\ ]$\\
\>Для каждого $u\in U[k,s]$:\\
\>\>Для каждого $tail\in$ Восстановить$(k+1,s-u)$:\\
\>\>\>Добавить $(u,tail)$ в $plans$\\
\>Вернуть $plans$\\
Ответ: $F[1,R]$ и Восстановить$(1,R)$
\end{tabbing}

\head{8. Программная реализация и диаграмма классов}
Самостоятельная реализация на Python без дополнительных библиотек.
Один класс хранит данные, таблицы значений и все оптимальные управления.
В коде этапам 1, 2, 3 соответствуют индексы 0, 1, 2.
\begin{center}\begin{tabular}{|p{0.89\linewidth}|}\hline
\centering\textbf{ResourceAllocation}\tabularnewline\hline
total: int; profits: list[list[int]]; lower, upper: list[int]\newline
n: int; values: list[list[number]]; choices: list[list[list[int]]]\tabularnewline\hline
\texttt{\_\_init\_\_(total, profits, lower, upper)}\newline
\texttt{solve()} $\to$ number\newline
\texttt{restore(k=0, s=None)} $\to$ list[tuple]\tabularnewline\hline
\end{tabular}\end{center}
\texttt{solve()} выполняет обратный проход; \texttt{restore()} восстанавливает
все планы. Внешние функции \texttt{show\_result(problem)} и \texttt{main()}
отвечают за вывод и запуск и не являются методами класса.
Поля \texttt{values} и \texttt{choices} заполняются в \texttt{solve()}.

\newpage
\head{9. Запуск, результат и ссылки}
Из каталога lab-04: \texttt{python python/dynamic\_programming.py}.
Ноутбук содержит тот же класс, отдельный блок данных и сохранённый вывод.
В Colab: «Среда выполнения $\to$ Выполнить все».
\begin{center}\begin{tabular}{c|rrrrrrrr}
$s$&0&1&2&3&4&5&6&7\\\hline
$F_3$&0&4&9&13&$-\infty$&$-\infty$&$-\infty$&$-\infty$\\
$F_2$&0&4&9&13&16&21&25&$-\infty$\\
$F_1$&$-\infty$&6&10&15&19&23&27&31
\end{tabular}\end{center}
Фактический запуск дал максимум 31 и планы $(1,3,3)$, $(2,3,2)$.
Вывод сохранён в \texttt{python/result.txt}. Независимый перебор шести
допустимых распределений подтвердил максимум и полноту списка.
Перебор используется только для проверки:
\begin{center}\begin{tabular}{c|r}
Распределение&Производительность\\\hline
$(1,3,3)$&31\\$(2,2,3)$&30\\$(2,3,2)$&31\\
$(3,1,3)$&28\\$(3,2,2)$&29\\$(3,3,1)$&29
\end{tabular}\end{center}
\href{https://github.com/Doney888/methods-of-optimization/tree/main/lab-04}{Репозиторий: lab-04, код и результат запуска.}

\href{https://colab.research.google.com/github/Doney888/methods-of-optimization/blob/main/lab-04/python/lab04-dynamic-programming.ipynb}{Открыть выполненный ноутбук в Google Colab.}

\href{https://drive.google.com/file/d/1oUGdd_os__GuU4qecMRgTcuhiaJS0kQc/view}{Исходное задание: Тема №4, задача 4.}

\head{10. Рефлективный вывод}
В работе удалось связать задачу распределения ресурса с состояниями и
переходами динамического программирования. Основная трудность --- учитывать
точное распределение всего ресурса и недостижимые остатки.
Значение $-\infty$ исключает такие продолжения из максимизации.
Ещё один важный момент --- сохранять все равные максимумы: иначе можно
потерять вторую оптимальную стратегию. Таблицы ручного решения и
программный расчёт совпали; перебор допустимых планов подтвердил полноту ответа.
\end{document}
